% Gerado de tcc.md (exportado do Claude Docs) por ferramentas/gera_latex.php em 12/11/2026 21:45.
% Para mudar o texto, edite este arquivo; se gerar de novo a partir do documento, as edições feitas aqui se perdem.

\chapter{FUNDAMENTAÇÃO TEÓRICA}\label{cap:2}

Este capítulo reúne os conceitos que sustentam o UNIA. As seções~\ref{sec:2.1} a~\ref{sec:2.3} tratam do direito à educação inclusiva, do PEI e dos níveis de suporte; as seções~\ref{sec:2.4} e~\ref{sec:2.5}, da medição do progresso, dos jogos digitais e dos trabalhos relacionados.

\section{Educação especial na perspectiva inclusiva}\label{sec:2.1}

A Constituição Federal garante o atendimento educacional especializado às pessoas com deficiência, preferencialmente na rede regular de ensino (BRASIL, [2026], art. 208, III). A Lei de Diretrizes e Bases da Educação Nacional (LDB) define a educação especial como modalidade escolar oferecida preferencialmente na rede regular e obriga os sistemas de ensino a assegurar currículos, métodos, técnicas e recursos educativos específicos para atender às necessidades desses educandos (BRASIL, 1996, arts. 58 e 59).

A Política Nacional de Educação Especial na Perspectiva da Educação Inclusiva, de 2008, deslocou o foco da escola especial para a classe comum: o AEE passou a complementar ou suplementar a escolarização, e não a substituí-la (BRASIL, 2008). O Decreto nº 7.611/2011 descreve o AEE como um conjunto de atividades e de recursos pedagógicos e de acessibilidade organizados de forma institucional e contínua (BRASIL, 2011), e a Resolução CNE/CEB nº 4/2009 atribui ao professor do AEE a elaboração, a execução e a avaliação do plano de AEE (BRASIL, 2009, art. 13).

A Lei Brasileira de Inclusão (LBI) ampliou esse quadro. Ela assegura um sistema educacional inclusivo em todos os níveis (BRASIL, 2015, art. 27) e inclui entre os deveres do poder público o planejamento de estudo de caso, a elaboração de plano de atendimento educacional especializado e a usabilidade pedagógica de recursos de tecnologia assistiva (BRASIL, 2015, art. 28, VII). A mesma lei define tecnologia assistiva como produtos, recursos, metodologias, estratégias, práticas e serviços que promovem a funcionalidade da pessoa com deficiência, com vistas à sua autonomia e inclusão (BRASIL, 2015, art. 3º, III).

Entre 2025 e 2026, três normas reorganizaram esse campo. O Decreto nº 12.686/2025 instituiu a Política Nacional de Educação Especial Inclusiva (PNEEI), voltada a estudantes com deficiência, com transtorno do espectro autista (TEA) e com altas habilidades ou superdotação (BRASIL, 2025a, art. 1º). O decreto adota o estudo de caso como “metodologia de produção, sistematização e registro de informações” sobre o estudante (BRASIL, 2025a, art. 11) e determina que a oferta do AEE não seja condicionada a diagnóstico ou laudo (BRASIL, 2025a, art. 11, § 7º). O Decreto nº 12.773/2025 alterou o texto para que o resultado do estudo de caso fundamente tanto o PAEE quanto o PEI, documentos obrigatórios e de atualização contínua (BRASIL, 2025b). Por fim, a Portaria MEC nº 421/2026 regulamentou a política, definiu o conteúdo mínimo de cada documento e determinou que o PAEE e o PEI sejam “revisados anualmente, de modo a compatibilizá-los com a avaliação contínua do estudante” (BRASIL, 2026, art. 7º, § 3º).

Três pontos dessas normas orientam este trabalho: o PEI deriva do estudo de caso, deve ser atualizado de forma contínua e precisa prever como o plano será acompanhado. O capítulo~\ref{cap:4} mostra como o UNIA trata cada um deles.

\section{Plano Educacional Individualizado}\label{sec:2.2}

O PEI é o registro do planejamento feito para um estudante específico. Tannús-Valadão e Mendes (2018) o descrevem como, ao mesmo tempo, um documento que especifica o programa educacional de um estudante a partir de seus pontos fortes e necessidades, um registro das diferenciações de que ele precisa para alcançar as expectativas de aprendizagem e um instrumento de prestação de contas ao aluno, à família e à equipe. Não se trata de um currículo paralelo: o plano parte do currículo comum e indica os apoios de que o aluno precisa para participar dele.

Francisco (2023) organiza o planejamento educacional individualizado em cinco etapas, que se repetem em ciclo:

\begin{passos}
  \item Identificação do estudante e de suas necessidades.
  \item Avaliação das habilidades que o estudante já realiza.
  \item Estruturação do plano, com as habilidades priorizadas e as estratégias.
  \item Aplicação das atividades planejadas.
  \item Avaliação e replanejamento, com a revisão do plano a partir dos resultados.
\end{passos}

As duas últimas etapas são as que mais dependem de registros. É nelas que se verifica se a habilidade priorizada avançou e se a estratégia deve mudar. Tannús-Valadão e Mendes (2018, p. 13) observam que estudos feitos na França, na Espanha, na Itália e nos Estados Unidos indicam efeitos positivos do uso do PEI e, principalmente, de seus critérios de acompanhamento, como a avaliação e a reavaliação do processo de escolarização.

A literatura aponta, porém, que essas etapas nem sempre se completam. Em 2018, Tannús-Valadão e Mendes constatavam que o Brasil não tinha dispositivos na legislação federal que garantissem um PEI baseado nas peculiaridades de cada estudante e que, na prática, o planejamento se concentrava nos serviços da sala de recursos multifuncionais, sem relação com o ensino na classe comum. As normas de 2025 e 2026, descritas na seção~\ref{sec:2.1}, passaram a exigir o PEI, mas a exigência, sozinha, não garante que o plano seja acompanhado. Um plano que não é revisto com base em dados tende a virar documento de arquivo.

A Portaria MEC nº 421/2026 deu forma a esse instrumento. Pelo art. 11, o PEI “contempla o plano de acessibilização curricular” e deve conter, no mínimo (BRASIL, 2026, art. 11):

\begin{alineasIFTO}
  \item as atividades a serem desenvolvidas na sala de recursos multifuncionais ou em outros espaços do AEE e sua articulação com o professor regente;
  \item as medidas de acessibilidade curricular, didático-pedagógica e avaliativa indicadas pelo estudo de caso;
  \item as estratégias de acompanhamento e de monitoramento do plano;
  \item o registro das devolutivas às famílias.
\end{alineasIFTO}

A norma exige estratégias de acompanhamento e de monitoramento, mas não diz como executá-las. É nesse ponto que um sistema pode ajudar, ao transformar a aplicação das atividades em registros que se acumulam e podem ser comparados.

\section{Níveis de suporte e autonomia}\label{sec:2.3}

Para registrar o que o aluno já realiza, o AppPEI usa um inventário de habilidades em que cada habilidade recebe um de três níveis: “não realiza”, “realiza com ajuda” e “realiza sem necessidade de suporte” (FRANCISCO, 2023). Os níveis dizem mais do que acerto ou erro: indicam quanto apoio o aluno precisa para chegar ao resultado.

Essa gradação tem base na teoria histórico-cultural. Vigotski (2007) distingue o nível de desenvolvimento real, expresso no que a criança resolve sozinha, do nível de desenvolvimento potencial, expresso no que ela resolve com a orientação de um adulto ou em colaboração com um colega mais capaz. A distância entre os dois é a zona de desenvolvimento proximal, e o que hoje a criança faz com ajuda tende a ser feito por ela sozinha mais tarde. Nessa leitura, a habilidade em “realiza com ajuda” é justamente aquela em que o ensino deve atuar, e a passagem para “realiza sem necessidade de suporte” é o sinal de que o apoio pode ser retirado.

No UNIA, a ajuda oferecida na tela é a dica: quando o aluno a pede, o sistema elimina opções e deixa apenas duas, uma delas a correta. Por ser um apoio padronizado, a dica pode ser contada, e a frequência com que o aluno a usa ao longo das sessões indica o quanto ele ainda depende de suporte. Trata-se de uma simplificação: em sala, a ajuda assume várias formas (verbal, gestual, física), e o sistema registra apenas a que ele próprio oferece. A regra que transforma acertos e dicas em níveis está na seção~\ref{sec:4.5}.

\section{Medição frequente do progresso e jogos digitais}\label{sec:2.4}

Saber se uma meta avança exige medir a mesma habilidade várias vezes, da mesma forma, e comparar os resultados. Deno (1985) propôs essa lógica na medida baseada no currículo (\textit{curriculum-based measurement}): provas curtas, aplicadas com frequência sobre as habilidades que estão sendo ensinadas, cujos resultados são lançados em gráfico para que o professor veja a tendência e decida se mantém ou muda o ensino. A proposta surgiu na educação especial como alternativa intermediária entre os testes padronizados e a observação informal, reunindo vantagens de ambos (DENO, 1985).

O UNIA segue essa lógica em escala reduzida. Cada sessão tem cinco rodadas sobre uma habilidade da meta, o resultado é gravado com data e as sessões de uma meta formam uma série que o sistema transforma em indicadores e gráficos. A diferença é que, no UNIA, quem aplica e corrige a prova é o próprio sistema.

Os jogos digitais são um meio adequado para isso. Prensky (2012) destaca que os jogos combinam envolvimento com retorno imediato: o jogador age, vê o resultado e tenta de novo. Para a medição, interessa uma propriedade adicional. Numa atividade digital, cada resposta pode ser registrada no momento em que acontece, com o acerto, o tempo e o apoio usado, sem trabalho extra do professor; no papel, o mesmo registro exigiria anotar cada tentativa durante a atividade.

Há dois limites a considerar. O primeiro é que o jogo mede o desempenho naquele jogo, e a transferência da habilidade para outras situações ainda precisa ser observada pelo professor. O segundo é que todo jogo depende dos sentidos e dos movimentos que exige: uma atividade que pede para ver figuras e tocá-las não mede da mesma forma o desempenho de uma criança com baixa visão ou com dificuldade motora. Esses limites voltam no capítulo~\ref{cap:5}.

\section{Trabalhos relacionados}\label{sec:2.5}

Os trabalhos relacionados foram localizados por uma busca estruturada, feita em 9 de novembro de 2026 e guiada por uma pergunta: que recursos digitais apoiam o PEI ou o ensino de estudantes com deficiência ou TEA na educação básica, e esses recursos registram o desempenho do estudante e o ligam às metas do plano? Foram consultadas três bases, sem limite de ano:

\begin{alineasIFTO}
  \item Biblioteca Digital Brasileira de Teses e Dissertações (BDTD), pela busca da própria base, em todos os campos: “plano educacional individualizado” aplicativo (19 resultados), “plano educacional individualizado” sistema digital (6) e “jogos digitais” “educação especial” alfabetização (6);
  \item SBC OpenLib, biblioteca digital da Sociedade Brasileira de Computação (SBC), por buscador da web restrito aos endereços da base: plano educacional individualizado sistema aplicativo (9 trabalhos) e jogo digital alfabetização crianças deficiência intelectual autismo desempenho professor acompanhamento (9);
  \item Scientific Electronic Library Online (SciELO), por buscador da web restrito ao endereço da base: plano educacional individualizado tecnologia digital aplicativo (7) e jogos digitais educação especial alfabetização matemática avaliação desempenho (5).
\end{alineasIFTO}

Descartados os resultados repetidos e os de fora das bases, restaram 57 trabalhos: 27 da BDTD, 18 da SBC OpenLib e 12 da SciELO. Leram-se o título e o resumo de cada um e, quando o resumo não bastou, o texto completo. Foram incluídos os trabalhos que desenvolvem um recurso digital (sistema, aplicativo ou jogo) voltado à educação básica e destinado a estudantes com deficiência ou TEA, ou à elaboração e ao acompanhamento do PEI. Foram excluídos os estudos sem recurso digital, que eram a maioria e tratavam de formação docente, legislação e prática pedagógica; os que avaliam jogos prontos, sem desenvolvê-los; os que oferecem material de apoio em vez de um recurso usado com o estudante ou com o PEI, como um e-book de estratégias, um repositório de orientações para professores e um modelo para quem desenvolve jogos; os voltados ao ensino superior, a rotinas diárias ou a estudantes fora da educação especial; os que não informam o público ou o conteúdo trabalhado; e um artigo cujo texto estava indisponível no dia da busca. Dez trabalhos atenderam aos critérios, cinco da BDTD e cinco da SBC OpenLib. Nenhum veio da SciELO, cujos resultados discutiam o PEI e a prática docente ou usavam jogos prontos ou não digitais. A eles se somaram dois trabalhos conhecidos antes da busca: o ABC Autismo, que atende aos mesmos critérios, e o sistema descrito na revista Temática, que gere os PEIs dos estudantes do IFRN, inclusive os dos cursos técnicos integrados ao ensino médio. Os doze trabalhos formam dois grupos: ferramentas para o PEI, a começar pelo AppPEI, e jogos ou aplicativos usados pelo estudante.

Priscila Rodrigues Cruz Francisco desenvolveu o AppPEI no Mestrado Profissional em Ensino na Educação Básica do Instituto de Aplicação Fernando Rodrigues da Silveira (CAp-UERJ). O recurso foi construído em planilha eletrônica do Google e tem nove menus: capa, informações, menu inicial, alunos, professores, avaliação diagnóstica, matriz e priorização de habilidades, proposta e PEI completo. A matriz de habilidades, construída pela autora a partir de referências da área, divide-se em quatro áreas: Alfabetização Matemática, Comunicação Oral, Leitura e Escrita e Alfabetização Digital. Cinco professores do colégio participaram da avaliação; quatro responderam ao questionário e três concluíram um PEI. Metade considerou o acesso fácil, 75\% consideraram o recurso intuitivo e a nota média foi 8,25. As sugestões incluíram migrar para páginas em HyperText Markup Language (HTML), melhorar a navegação e o layout, permitir o uso no celular e deixar o professor editar a matriz de habilidades (FRANCISCO, 2023).

Dois trabalhos da busca também tratam o PEI como documento digital. Almeida (2023) acompanhou professores do ensino comum e da educação especial na elaboração do PEI de um aluno autista, numa escola de Balsas, no Maranhão, e propôs como produto o PEI-DIGITAL, aplicativo de uso interno da escola que os professores podem preencher em conjunto depois de definir os objetivos, as metas, as estratégias e as adaptações curriculares do plano. Diniz \textit{et al.} (2024, p. 107-109) apresentaram o PAR Digital, sigla de Planejar, Aplicar, Rever: um sistema web para preencher e gerir o PEI de estudantes com deficiência em classes comuns, com um banco de estratégias baseadas no Desenho Universal para a Aprendizagem. O professor cadastra as estratégias do plano e marca a situação de cada uma, como alcançada ou não alcançada, e o sistema gera relatórios de progresso por disciplina; segundo os autores, a ferramenta foi validada e usada por escolas da educação básica.

Paiva \textit{et al.} (2026) desenvolveram, por pesquisa-ação e em parceria com o Núcleo de Apoio às Pessoas com Necessidades Educacionais Específicas (NAPNE) do Instituto Federal do Rio Grande do Norte (IFRN), Campus Pau dos Ferros, um sistema web para gerir os PEIs dos estudantes da instituição. O PEI é criado automaticamente quando o estudante é incluído na oferta de uma disciplina; o professor registra os objetivos, os conteúdos, as metodologias e os recursos adaptados, e o plano passa pelas situações não iniciado, em andamento, preenchido e finalizado. O acompanhamento é feito pelos pareceres pedagógicos de cada bimestre ou semestre, pelas notas, pelo histórico do estudante e por comentários da equipe (PAIVA \textit{et al.}, 2026, p. 65, 69-73). Na avaliação pelo Modelo de Aceitação de Tecnologia, os 10 profissionais participantes concordaram, total ou parcialmente, que aprender a usar o sistema é fácil (PAIVA \textit{et al.}, 2026, p. 78). O trabalho mostra que a demanda por apoio digital ao PEI também aparece na rede federal.

O segundo grupo reúne jogos e aplicativos usados pelo estudante. Farias, Silva e Cunha (2014), do Instituto Federal de Alagoas (IFAL), desenvolveram o ABC Autismo, aplicativo móvel que apoia a alfabetização de crianças com autismo com base nas premissas do programa Treatment and Education of Autistic and Related Communication Handicapped Children (TEACCH); a validação com crianças, pais e profissionais de uma associação indicou alinhamento ao programa e interface autoexplicativa. Negrim (2019) desenvolveu um jogo digital de alfabetização e o usou com cinco alunos com deficiência intelectual do ensino fundamental; avaliações diagnósticas feitas ao longo das fases do jogo apontaram avanços nas hipóteses de escrita da maioria deles. Bentes, Araújo e Pires (2020) criaram o ANE, aplicativo de jogos para que crianças com deficiência intelectual aprendam o alfabeto, os numerais e as emoções, com requisitos levantados a partir de dados do atendimento educacional especializado de um município; nos testes com usuários finais, todos o consideraram boa ferramenta de apoio para responsáveis e profissionais da educação especial. Dutra, Maschio e Gasparini (2023) desenvolveram o Pensar e Lavar, jogo acessível de pensamento computacional para crianças neurotípicas e com deficiência intelectual no início da alfabetização, avaliado em quatro etapas, uma delas com 47 crianças.

Galdino (2024, p. 61, 66) criou um site com jogos digitais de alfabetização matemática para alunos com TEA, como os de sequência numérica e de quantidade, feitos no PowerPoint para que o professor possa editá-los e adaptá-los a cada aluno. Silva, Santos e Cunha (2024), também do IFAL, apresentaram o LetraMundo, jogo sério web para a alfabetização de crianças com TEA em ambiente escolar; no teste preliminar, profissionais da educação especial apontaram a necessidade de redistribuir os níveis de dificuldade e de criar atividades com grafemas. Mesquita e Bastos (2025) descreveram o Garimpeirinhos, ambiente web de apoio à alfabetização e à inclusão na educação infantil e nos anos iniciais, com jogos, leitura progressiva, prancha de comunicação alternativa e utilitários para o professor. Fernandes (2026) desenvolveu cinco jogos no Scratch para a alfabetização de estudantes com TEA, com foco em consciência fonológica, relação entre grafema e fonema e sequência alfabética; a avaliação, qualitativa, apontou mais engajamento e mais tempo de atenção durante as tarefas.

\begin{tabelaIFTO}
\setlength{\largtab}{\dimexpr\linewidth-10\tabcolsep-6\arrayrulewidth-1pt\relax}
\begin{longtable}{|L{0.2253}|L{0.16}|L{0.2253}|L{0.2008}|L{0.1885}|}
\legendatab{quadro}{Trabalhos relacionados e o UNIA}{qua:1}{5}\\
\hline
\textbf{Trabalho} & \textbf{Recurso} & \textbf{Público e foco} & \textbf{Registro do desempenho} & \textbf{Ligação com as metas do PEI} \\
\hline
\endfirsthead
\multicolumn{5}{@{}r@{}}{\textit{(continuação)}}\\
\hline
\textbf{Trabalho} & \textbf{Recurso} & \textbf{Público e foco} & \textbf{Registro do desempenho} & \textbf{Ligação com as metas do PEI} \\
\hline
\endhead
\multicolumn{5}{@{}r@{}}{\textit{(continua)}}\\
\endfoot
\fontetab{5}{elaboração própria (2026), com base nos trabalhos citados. “Não relatado” e “Não abordada” indicam que o tema não aparece no texto consultado de cada trabalho.}
\endlastfoot
AppPEI (FRANCISCO, 2023) & Planilha eletrônica & Professores; elaboração do PEI & Avaliação diagnóstica preenchida pelo professor & Monta o PEI a partir das habilidades priorizadas \\ \hline
PEI-DIGITAL (ALMEIDA, 2023) & Aplicativo & Professores; elaboração do PEI em conjunto & Não relatado & Registra objetivos, metas e estratégias do plano \\ \hline
PAR Digital (DINIZ \textit{et al.}, 2024) & Sistema web & Professores e coordenadores; gestão do PEI & Situação de cada estratégia, marcada pelo professor, e relatórios & Estratégias e relatórios ficam no PEI \\ \hline
Sistema de gestão de PEIs do NAPNE/IFRN (PAIVA \textit{et al.}, 2026) & Sistema web & Professores, pedagogos e coordenação; elaboração e gestão do PEI por disciplina & Pareceres por bimestre ou semestre e notas, lançados pela equipe & Objetivos, conteúdos e métodos adaptados ficam no PEI; o progresso não é medido por meta \\ \hline
ABC Autismo (FARIAS; SILVA; CUNHA, 2014) & Aplicativo móvel & Crianças com TEA; alfabetização & Não relatado & Não abordada \\ \hline
Jogo de alfabetização (NEGRIM, 2019) & Jogo digital & Alunos com deficiência intelectual; leitura e escrita & Avaliações diagnósticas ao longo das fases do jogo & Não abordada \\ \hline
ANE (BENTES; ARAÚJO; PIRES, 2020) & Aplicativo de jogos & Crianças com deficiência intelectual; alfabeto, numerais e emoções & Não relatado & Não abordada \\ \hline
Pensar e Lavar (DUTRA; MASCHIO; GASPARINI, 2023) & Jogo digital & Crianças com e sem deficiência intelectual; pensamento computacional & Não relatado & Não abordada \\ \hline
Jogos de alfabetização matemática (GALDINO, 2024) & Site com jogos feitos no PowerPoint & Alunos com TEA; alfabetização matemática & Não relatado & Discutida no texto, não nos jogos \\ \hline
LetraMundo (SILVA; SANTOS; CUNHA, 2024) & Jogo web & Crianças com TEA; alfabetização na escola & Não relatado & Não abordada \\ \hline
Garimpeirinhos (MESQUITA; BASTOS, 2025) & Ambiente web & Educação infantil e anos iniciais, com recursos para TEA; alfabetização & Não relatado & Não abordada \\ \hline
Jogos em Scratch (FERNANDES, 2026) & Jogos digitais & Estudantes com TEA; consciência fonológica e alfabeto & Observação e registros das professoras & Não abordada \\ \hline
UNIA (este trabalho) & Sistema web & Professor e aluno dos anos iniciais; metas do PEI & Acertos, erros, dicas e tempo de cada sessão, gravados pelo sistema & Cada sessão medida entra no histórico da meta \\ \hline
\end{longtable}
\end{tabelaIFTO}

O Quadro~\ref{qua:1} mostra dois grupos. As ferramentas do primeiro organizam o PEI como documento. A mais próxima do UNIA é o PAR Digital, que liga as estratégias ao plano e gera relatórios de progresso, mas o andamento depende do que o professor marca em cada estratégia (DINIZ \textit{et al.}, 2024, p. 109). No sistema do NAPNE/IFRN, o progresso também depende do que a equipe registra: os pareceres de cada bimestre ou semestre e as notas (PAIVA \textit{et al.}, 2026, p. 70, 72). Os jogos do segundo grupo trabalham as habilidades com o estudante; nos que relatam algum acompanhamento, ele vem de avaliações da pesquisa ou da observação das professoras, e nenhum liga o desempenho às metas de um plano individualizado. Fernandes (2026) recomenda, como próximo passo, incorporar métricas padronizadas de aprendizagem e registrar de forma sistemática o progresso dos estudantes. O UNIA junta as duas partes: os mini games existem para medir as metas do PEI, o próprio sistema grava o desempenho do aluno em cada sessão, e cada resultado entra no histórico da meta correspondente.
